\section{Original-policy action comparisons and their information cost}
\label{supp:improvement120}

This section proves the R120 comparison construction under the original
investment primitives.  The completed-incumbent safeguard, three projected
support proposals, and original all-action witnesses retain their
R108--R110 ownership.  The common-risk identity retains its R111/R119
ownership.  The new ingredients proved here are an original AVaR
comparison modulus, the associated mixed integration interface and its
restricted information bound, and quantitative control of common-tail
mismatch.  The statistical transfer and strong-convexity safeguards use
classical arguments and are proved to make their assumptions explicit.
No statement supplies a neural representation separation.

\subsection{Valid paired integration and the policy contrast}
\begin{proof}[Proof of Theorem~\ref{thm:paired120}]
Write $D=f_b-f_a$.  A Lipschitz function of constant $E_{ab}$
satisfies on a cell with midpoint $z_i$ and length $h_i$
\[
 \left|\int_{B_i}D(z)\,dz-h_iD(z_i)\right|
 \le E_{ab}\int_{B_i}|z-z_i|\,dz
 =E_{ab}h_i^2/4.
\]
Summing and dividing by $B=2Rt$ proves the second interval in
\eqref{eq:pairintersection120}, including arbitrary valid midpoint
value errors.  Separately, convexity gives for each $f_a$
\[
 h_if_a(z_i)\le\int_{B_i}f_a(z)\,dz
       \le\frac{h_i}{2}\{f_a(l_i)+f_a(r_i)\}.
\]
Valid lower midpoint and upper endpoint values thus give
$L_a\le\mathbb E_Pq_uf_a\le U_a$, and likewise for $b$.
Subtracting these two valid intervals proves the first interval.
Both contain the same integral, so their exact intersection is
nonempty and valid.  Empty outward intersections signal a failed
numerical construction, not a license to choose a favorable endpoint.
Zero-length tails contribute zero measure and are omitted.

For completeness, each exact convex midpoint--secant width is
bounded by $KRt/n^2$ under $h_i\le B/n$.  On a cell its unnormalized
gap is $h_i^2/4$ times the difference of the right and left
half-cell secant slopes.  These differences are nonnegative; ordering
all cells, including across the omitted center, bounds their sum by
$2K$.  Thus the normalized width is at most
$Kh_{\max}^2/(2B)\le KRt/n^2$.
Two such widths sum to $2KRt/n^2$.
The paired Lipschitz interval has width
\[
 \frac{E_{ab}}{2B}\sum_i h_i^2
 \le E_{ab}h_{\max}/2\le E_{ab}Rt/n.
\]
Intersection proves \eqref{eq:mixedwidth120}.
The midpoint--secant inequality is applied only to each convex
continuation, never to $D$.

For the original continuations, Theorem~\ref{thm:modulus120}
also gives a Lipschitz constant $2M_\alpha$ for $f_a'$:
the innovation direction has squared norm two.
The two Taylor upper bounds at the endpoints of a cell
therefore imply
\[
 \frac{f_a(l_i)+f_a(r_i)}2-f_a(z_i)
 \le \frac{2M_\alpha h_i^2}{8}.
\]
The opposite linear terms cancel.  Multiply by $h_i/B$
and sum over the cells.  The midpoint--secant width for
one function is at most $M_\alpha h_{\max}^2/4$,
and the difference of the two separately enclosed
integrals has width at most
$M_\alpha h_{\max}^2/2\le2M_\alpha R^2t^2/n^2$.
Intersecting this analytic bound with the previous
ones proves \eqref{eq:smoothpairwidth120}.
The numerical algorithm already computes these same
midpoint and endpoint values; this consequence changes
neither its inputs nor its prescribed decisions.

For the same feasible original density, define
$k_a=\rho_\alpha(f_a)-\mathbb E_Pq_uf_a\ge0$ and similarly $k_b$.
Direct cancellation yields
\[
 \rho_\alpha(f_b)-\rho_\alpha(f_a)
     =\mathbb E_Pq_uD+k_b-k_a.
\]
Consequently the optimal-future root contrast lies in
$h(b)-h(a)+\beta[d_--\kappa_a,d_++\kappa_b]$.
By \eqref{eq:future120}, the difference of the two actual future
allowances lies in $[-\delta,\delta]$, not $[-2\delta,2\delta]$.
Adding this interval proves \eqref{eq:policycontrast120}.
This inference uses the same implemented future policy at both
actions and makes no convexity assertion about its value.
When the two realized actions agree, their complete-policy values
and their continuation functions agree exactly.  Their difference
is zero before any interval widening or new acquisition.
\end{proof}

\paragraph{Explicit node-precision allowance.}
Suppose every acquired scalar value interval has width at most
$\xi$ and the sums are computed exactly.  The paired midpoint
interval gains width at most $2\xi$ from its node intervals.
The separate convex difference gains width at most $4\xi$.
Thus a valid bound replacing \eqref{eq:mixedwidth120} is
\begin{equation}\label{eq:inexactpairwidth120}
 \min\left\{\frac{2KRt}{n^2}+4\xi,
                      \frac{ERt}{n}+2\xi\right\}.
\end{equation}
To prove this, write each midpoint lower error and each endpoint
upper error separately.  Their nonnegative normalized weights sum
to one, so each such weighted error is at most $\xi$.
For the paired midpoint difference the sum of the two interval
widths is at most $2\xi$.  The remaining terms are the exact
quadrature gaps just proved.  Outward accumulation adds its
separately bounded arithmetic enlargement.  In an implementation,
the directly retained interval endpoints remain authoritative;
neither $\xi$ nor the future-policy allowance is silently sent to zero.

\subsection{A matching bound for the fixed-tail information interface}
\begin{proposition}[Paid paired point information]
\label{prop:pairinformation120}
Fix $R,K>0$, $t=1-\alpha>0$, $\alpha\in[0,1)$, and $E\ge0$.
The unknown pair consists of convex $K$-Lipschitz functions on
$[-R,R]$ whose difference is $E$-Lipschitz.  A known feasible
two-tail density and these class restrictions are supplied.
The information oracle returns an exact value of one chosen
function at one chosen point.  A reported interval must contain
the fixed-density integral of the difference for every pair
consistent with its transcript.

For every integer $n\ge1$, there is a nonadaptive quadrature
algorithm using at most $4n+8$ scalar values whose width is at most
$\min\{2KRt/n^2,ERt/n\}$.
Conversely, for the known one-sided density
$q=t^{-1}\boldsymbol1_{[-R+2R\alpha,R]}$, every deterministic
adaptive algorithm using at most $N$ scalar values has an input
requiring width at least \eqref{eq:pairlower120}.
For $N\ge32$ the upper width can in particular be bounded by
\[
 \min\{128KRt/N^2,\ 8ERt/N\}.
\]
These are interval-width bounds, not point-estimation error bounds.
\end{proposition}

\begin{proof}
For the upper bound, divide a nonempty tail of length $b_j$ into
$\lceil nb_j/(2Rt)\rceil$ equal cells.  The two tail lengths sum
to $2Rt$, so at most $n+1$ nonempty cells suffice.  Each function
requires at most $2n+4$ midpoint and distinct endpoint values.
Paying for both functions gives $4n+8$; Theorem~\ref{thm:paired120}
proves its width.  At $\alpha=0$ this is just integration over the
whole interval and requires no tail location.  If $N\ge32$, take
$n=\lfloor(N-8)/4\rfloor\ge N/8$.  The displayed upper bound
follows.

For the lower bound put $q_0=-R+2R\alpha$ and $B=R-q_0=2Rt$.
Run the proposed algorithm on the pair $(f_0,f_0)$, where
\[
 f_0(z)=\frac{K}{2B}(z-q_0)_+^2.
\]
It is convex, nondecreasing, and $K$-Lipschitz.  Partition its
right tail into $2(N+1)$ equal cells of length
$h=B/\{2(N+1)\}$.  At most $N$ cells contain a queried location
in their interior, irrespective of which function label was
queried.  Choose $N+1$ cells containing no such point.  On each
selected cell replace $f_0$ by the convex combination
\[
 f_1=(1-\lambda)f_0+\lambda\,(\text{chord of }f_0),
 \qquad
 \lambda=\min\{1,4E(N+1)/K\},
\]
and put $f_1=f_0$ elsewhere.
Every queried value remains unchanged, including queries at
selected cell endpoints.  Induction on the adaptive transcript
therefore gives the same queries and answers on $(f_0,f_1)$.

The chord slope is between the two endpoint derivatives of
$f_0$.  At entry to a selected cell the derivative of $f_1$
jumps upwards by $\lambda Kh/(2B)$; at its exit it jumps
upwards by the same amount.  Interior derivatives are
nondecreasing, and all slopes lie in $[0,K]$.  Adjacent selected
cells have the same ordering.  Hence $f_1$ is convex and
$K$-Lipschitz.  On a selected cell
\[
 |(f_1-f_0)'|\le\lambda Kh/(2B)
             =\lambda K/\{4(N+1)\}\le E.
\]
The difference is continuous and zero at every selected
endpoint and elsewhere, so it is globally $E$-Lipschitz.
The construction thus preserves every information restriction.

The integral of the chord excess of a quadratic with second
derivative $K/B$ on a cell of length $h$ is $Kh^3/(12B)$.
Summing $N+1$ selected cells and dividing by $B$ gives
\begin{align*}
 \mathbb E_Pq(f_1-f_0)
 &\ge\frac{\lambda KB}{96(N+1)^2}\\
 &=\min\left\{\frac{KRt}{48(N+1)^2},
                       \frac{ERt}{12(N+1)}\right\}.
\end{align*}
The first pair has contrast zero.  A sound interval for their
common transcript must contain both contrasts and hence has
at least this width.  Both functions are nondecreasing, so
the known right tail is also exactly maximizing for both
AVaRs.  This supplies a subclass with no localization slack;
it does not eliminate localization costs from the upper
algorithm on a general pair.
\end{proof}

The proof uses the standard indistinguishable-input approach to
numerical information \citep{Novak2015Integration97}.  Its upper
bound is conditional on a known fixed tail, and its lower bound
allows arbitrary convex functions with the stated moduli.  It
does not claim that the original finite-parameter investment
primitives realize this entire adversarial class.  Formula access,
derivative access, additional separate smoothness restrictions,
a shared exact symbolic representation, and free terminal
solves are different information conventions.
Providing one of them changes the comparison for every method,
including splines, polynomials, hinges, and neural networks.
In particular \eqref{eq:smoothpairwidth120} uses an additional
proved property of the original continuation.  No matching
lower bound for that full original-primitives information
class is inferred from the nonsmooth construction above.

\subsection{Transferring the defect of a common tail}
\begin{theorem}[Action-scaled transfer between constrained tail witnesses]
\label{thm:tailtransfer120}
Let $0<\alpha<1$, $L=2R\alpha$, $B=2Rt$, and
$I=[-R,R-L]$.  Suppose $f_a,f_b$ are $m_z$-strongly convex on
the innovation interval, with $m_z>0$, and their difference is
$E$-Lipschitz.  Define
$H_j(u)=\int_u^{u+L}f_j(z)\,dz$ and
$k_j(q_u)=\{H_j(u)-\min_IH_j\}/B$.
For any feasible $u$ and any verified $\kappa_b\ge k_b(q_u)$,
\[
 k_a(q_u)\le
 \left\{\sqrt{\kappa_b}
                  +E\sqrt{\frac{\alpha}{2m_zt}}\right\}^{2}.
\]
The statement includes an endpoint optimizer or a one-sided tail.
For the original investment pair, $m_z=181/64$ and
$E=E_{ab}$ are valid.  At $\alpha=0$ the only relevant
density is the full original law and both deficits are zero.
\end{theorem}

\begin{proof}
The original tail representation, or its direct proof in
Theorem~\ref{thm:tailinformation119}, gives the stated expression
for the risk deficit.  Strong convexity implies that $H_a$
and $H_b$ are $\lambda$-strongly convex on $I$, where
$\lambda=m_zL$.  Indeed
$H_j'(u)=f_j(u+L)-f_j(u)$, and its increment between two
cuts is at least $m_zL$ times their distance.
Let $u_b$ minimize $H_b$ on $I$.  The constrained first-order
inequality, followed by strong convexity, gives
\[
 H_b(v)\ge H_b(u_b)+\tfrac12\lambda|v-u_b|^2
 \quad(v\in I).
\]
There is no assumption that $H_b'(u_b)=0$ at a boundary.
Write $G=H_b-H_a$.  Since
$G'(v)=(f_b-f_a)(v+L)-(f_b-f_a)(v)$,
$G$ has Lipschitz constant $EL$.

Put $e=H_b(u)-H_b(u_b)=Bk_b(q_u)\le B\kappa_b$.
The preceding inequality gives
$|u-u_b|\le\sqrt{2e/\lambda}$.
For any $v\in I$, the triangle inequality now yields
\begin{align*}
 H_a(u)-H_a(v)
 &\le e-\tfrac12\lambda|v-u_b|^2+EL|v-u|\\
 &\le e+EL\sqrt{2e/\lambda}
           -\tfrac12\lambda|v-u_b|^2+EL|v-u_b|\\
 &\le e+EL\sqrt{2e/\lambda}+\frac{E^2L^2}{2\lambda}.
\end{align*}
The last step maximizes a scalar quadratic over a larger
nonnegative-distance domain, so it is valid at either boundary.
Take the supremum in $v$, divide by $B$, and use
$L/B=\alpha/t$.  The result is
$k_a(q_u)\le\{\sqrt{k_b(q_u)}
          +E\sqrt{\alpha/(2m_zt)}\}^2$.
Monotonicity in $k_b$ proves the asserted bound.
\end{proof}

\paragraph{Outward square roots in the transferred bound.}
The implementation intersects the directly acquired incumbent
deficit with an outward version of the transfer bound.
For an exact nonnegative rational $z=p/q$ with $q>0$,
let $m$ be the integer square root of
$\lfloor 2^{80}p/q\rfloor$, and increase $m$ by one if
$m^2q<2^{80}p$.  The resulting $m/2^{40}$ is at least
$\sqrt z$: before a possible increment, $m$ is the
floor of $2^{40}\sqrt z$, and the exact square comparison
detects whether that floor was exact.  A binary64 conversion
is moved upwards when its exact rational value is smaller
than the dyadic bound, and its squared exact rational value
is checked against $z$.  Applying this procedure to upper
endpoints, followed by outward additions, products, and
squaring, gives a valid transferred deficit.
Taking the smaller of two valid upper bounds remains valid.
Both direct endpoint payloads and all integer operations are
still paid.  Because this is an incumbent-deficit bound,
its tightening directly affects the lower contrast and
the certified-refusal predicate; it does not itself tighten
the candidate-deficit term in the upper acceptance endpoint.

\subsection{Uniform safeguards, finite refinement, and an honest stopping floor}
\begin{proof}[Proof of Proposition~\ref{prop:safe120}]
At any accepted step the previous theorem gives
$J_x(b)-J_x(a)\le\overline D<0$, and hence the original
cost decreases by at least $-\overline D>0$.
The replacement is explicitly feasible and qualified.
A refused trial changes no action or future policy.
Induct over the finite sequence, and telescope the original
cost differences.  This gives the lower gain account in
\eqref{eq:gainaccount120}.  The initial completed policy satisfies
$J_x(a_{\mathrm{initial}})-V_0(x)\le\tau$, while every final
feasible policy has cost at least $V_0(x)$.  Subtraction gives
the upper bound $\tau$.

This proof is statewise and representation independent.
Finite Borel proposal, arithmetic, ordering with fixed tie rules,
and acceptance maps compose into a Borel returned-action rule,
provided the completed incumbent service is defined statewise.
The original finite-horizon policy theorem then defines its
policy value.  Exhaustion of a new comparison returns the
already completed incumbent.  A failed scientific execution
does not license omitting that requested state from the sample
or presenting an incomplete service as a completed policy.
\end{proof}

\begin{proposition}[Finite precision above the full comparison margin]
\label{prop:resolution120}
Fix two distinct feasible actions and one acquired feasible cut.
Let $g_Q=Q_x(a)-Q_x(b)$ and suppose $\kappa_b$ is a valid,
fixed upper bound on its candidate deficit.  Put
\[
 \gamma=E_{ab}\sqrt{\frac{\alpha}{2m_zt}},\qquad
 F_q=\kappa_b+(\sqrt{\kappa_b}+\gamma)^2.
\]
Assume refinement at positive integer resolutions $n$ produces
fixed-density contrast intervals of width
\[
 w_n\le \eta_{\mathrm{num}}+\min\{A_2/n^2,A_1/n\},
 \quad A_2=\min\{2KRt,2M_\alpha R^2t^2\},
 \quad A_1=E_{ab}Rt,
\]
where $\eta_{\mathrm{num}}\ge0$ includes the stated terminal-node
and accumulation errors.  Suppose
\[
 \epsilon=(g_Q-\delta)/\beta-F_q-\eta_{\mathrm{num}}>0.
\]
Then the upper contrast gate is strictly negative whenever
\begin{equation}\label{eq:resolution120}
 n>\min\{\sqrt{A_2/\epsilon},\,A_1/\epsilon\}.
\end{equation}
Acceptance also requires the separate original candidate
qualification; if it is established by a resolution $n_{\rm qual}$
and remains valid at the considered round, the two conditions
together suffice.

If the actual work of each doubling round is at most $c n+c_0$,
with $c,c_0\ge0$ and a positive integer starting resolution
$n_0$, the total through a final round $n_J$ is at most
$c(2n_J-n_0)+c_0(J+1)\le2c n_J+c_0(J+1)$,
apart from separately paid setup.
A first doubling round meeting an integer target $n_*\ge n_0$
satisfies $n_J<2n_*$ and
$J\le\lceil\log_2(n_*/n_0)\rceil$.
This cost statement is conditional on the declared work model.
Actual terminal solves, proposals, refused rounds, recording,
and qualification costs are not replaced by a point-query count.
\end{proposition}

\begin{proof}
Let $d=\mathbb E_Pq_u(f_b-f_a)$ and denote the true deficits
by $k_a,k_b$.  If $d_+$ is the reported upper integral bound,
$0\le d_+-d\le w_n$.
The complete-policy upper contrast used by the gate is
$U=h(b)-h(a)+\beta(d_++\kappa_b)+\delta$.
The exact common-risk identity gives
\begin{align*}
 U
 &=-g_Q+\delta
       +\beta\{(d_+-d)+(\kappa_b-k_b)+k_a\}\\
 &\le-g_Q+\delta+\beta\{w_n+F_q\},
\end{align*}
where Theorem~\ref{thm:tailtransfer120} bounds $k_a$ and
$k_b\ge0$ was retained in the first line and dropped only
for the upper estimate.  Under the stated positive margin,
\eqref{eq:resolution120} makes the right side strictly
negative.  No inference about candidate qualification follows
from that sign alone; its separate verified predicates are
part of the assertion.

For $n_j=2^jn_0$,
$\sum_{j=0}^J n_j=2n_J-n_0<2n_J$.
Sum the per-round work bound.  Minimality of the first
doubling round meeting $n_*$ gives the two displayed bounds
on $n_J$ and $J$.  For the finite executed ladder
$2,4,8,16,32,64$, this guarantee applies only if a qualifying
round also lies within that declared cap.
\end{proof}

The fixed-cut theorem does not assume that increasing the cell count
also improves its tail localization.  If the cut and its deficit
remain fixed, their allowances remain in $F_q$.  The actual
implementation retains both direct original deficits and
intersects the incumbent bound with the transferred one.
A failed
sufficient margin can therefore be conservative; it is not proof
that a particular candidate is economically inferior.

\begin{proposition}[Sharp residual uncertainty of the unchanged-future interface]
\label{prop:futurefloor120}
Suppose the only information about the actual-policy correction at
two distinct actions is $0\le\eta_x\le\delta$ in
\eqref{eq:future120}, and their two $Q_x$ values are known
exactly.  The interval
$Q_x(b)-Q_x(a)+[-\delta,\delta]$ is then sharp over the
compatible corrections.  Thus $g_Q>\delta$ is necessary
and sufficient for this information alone to certify strict
improvement uniformly.  Further refinement of the $V_1$
integration alone cannot remove that floor.
\end{proposition}

\begin{proof}
The difference $\eta_x(b)-\eta_x(a)$ is in
$[-\delta,\delta]$.
The upper endpoint is attained in the interface by setting
$\eta_x(a)=0$, $\eta_x(b)=\delta$; reversing those values
attains the lower endpoint.  If continuity were additionally
required without any quantitative modulus, the two prescribed
values could still be joined by a continuous bounded
interpolation between these distinct actions.
Consequently a strict negative upper contrast is equivalent
to $g_Q>\delta$.  The claim is about compatibility with the
declared information: it does not assert that both extremal
corrections are generated by the actual fixed terminal policy.
That policy's further verified structure would be additional
information.

For exactly the same action, both values of $\eta_x$ are the
same and cancellation is exact; the two-endpoint construction
does not apply.  If only the magnitude of the actual gain
$g_J$ is specified, then $g_Q\ge g_J-\delta$.
The corresponding sufficient worst-case condition for the
exact-$Q$ gate is $g_J>2\delta$, rather than $g_J>\delta$.
This distinction prevents a known-continuation gain and an
actual-policy gain from being assigned the same resolution
claim.
\end{proof}

\section{Population gain with frozen decisions and finite numerical intervals}
\label{supp:population120}
\begin{proof}[Proof of Proposition~\ref{prop:population120}]
Let $\mathcal F$ contain the frozen actors, fitting data,
proposal and refinement protocol, representation choices,
initial-state sampling law, and statewise execution rule.
Condition on $\mathcal F$.
For each of the $48$ fixed groups, require that its completed
incumbent and final action be measurable functions of its
own initial state, with the uniformly safe guarantee in
Proposition~\ref{prop:safe120}.
In particular they are not fitted, selected, or changed
after observing the validation states.  Apply each complete
service to a single original state; batching innovations
within that state does not violate this condition.
Batch logic that lets another sampled state's failed
certificate alter the current state's action would
require a different probability argument.

For the declared finite dyadic law one may take
\[
 \mathcal X_{40}
 =\{0,1,\ldots,2^{40}-1\}^2/2^{40}.
\]
Independent uniform coordinate draws from this set make
$X_1,\ldots,X_{64}$ iid under the finite uniform population.
A fixed sequence of $640$ independently uniform bytes
can encode these $128$ coordinates in blocks of five bytes.
Their later retention is reproducibility evidence.
Independence and uniformity of the generating entropy, and
its independence of $\mathcal F$, remain the statistical
assumption.  The conclusion is for this declared law, not
automatically for a continuous uniform population or an
empirical economic state distribution.

Write $G_j(x)$ for the actual complete-policy gain over
the corresponding completed incumbent.
The uniform safeguard gives $0\le G_j\le\tau$ before
the sample is drawn.  Let
$\ell_{ji}\le G_j(X_i)\le u_{ji}$
be the saved valid gain bounds, intersected with
$[0,\tau]$.  For example, accepted-step lower gains can
be summed outwardly, and the two separately valid
complete-policy value intervals can provide another gain
interval to intersect.  Numerical interval widths are
retained, not replaced by midpoints.  Exact equal-policy
identities give $G_j=0$ without independent interval
subtraction.  A missing independent numerical validation
can conservatively use $[0,\tau]$ and must not be dropped.
The numerical enclosure procedure may depend on already
observed validation information, provided its intervals
remain valid and it does not change the frozen returned
policy rule.

For a fixed group and threshold define the latent events
$B_i=\boldsymbol1_{\{G_j(X_i)\ge\theta_k\}}$ and their
probability $p_{jk}$.
These are iid Bernoulli variables conditional on
$\mathcal F$.
The observable certified count satisfies
\[
 K_{\mathrm{cert}}
 =\sum_i\boldsymbol1_{\{\ell_{ji}\ge\theta_k\}}
 \le\sum_iB_i=:K_{\mathrm{true}}
\]
pointwise.  The certified count itself need not be
binomial if the numerical validation is incomplete
or adaptive.  Its domination by the latent binomial
count is sufficient.  The parameters may be dependent
across groups; no cross-group independence is required.

Here is the exact inversion used in the statement.
For an observed count $k\in\{1,\ldots,n\}$, $n=64$, define
\[
 T_k(p)=\sum_{r=k}^n {n\choose r}p^r(1-p)^{n-r}.
\]
It is continuous and strictly increasing on $(0,1)$, with
endpoint values zero and one.  Differentiating the finite
sum and telescoping proves
$T_k'(p)=n{n-1\choose k-1}p^{k-1}(1-p)^{n-k}>0$.
There is therefore a unique solution of
$T_k(p)=1/7680$ in $(0,1)$; call it $p_*(k)$.
For $k=0$ set $p_*(0)=0$.
The map $k\mapsto p_*(k)$ is nondecreasing because
$T_{k+1}(p)\le T_k(p)$.
For the true parameter $p$, the event
$p_*(K)>p$ is either empty or of the form $K\ge k_0$,
where $T_{k_0}(p)<1/7680$.
Its probability is therefore at most $1/7680$.

An exactly rational implementation need not evaluate a
beta quantile.  Start with the bracket $[0,1]$ and
bisect it $48$ times, evaluating the finite binomial
polynomial as an exact rational.  Move its lower endpoint
to the midpoint only when $T_k$ there is at most $1/7680$.
The returned lower endpoint $p^-$ satisfies
$p^-\le p_*(k)$ and $p_*(k)-p^-\le2^{-48}$.
Thus it retains the same conservative coverage when
applied to $K_{\mathrm{true}}$.
Monotonicity of $p_*(k)$ and
$K_{\mathrm{cert}}\le K_{\mathrm{true}}$ imply
$p^-(K_{\mathrm{cert}})\le p_*(K_{\mathrm{true}})$.
Consequently the same coverage holds for the actually
reported certified count, including any conservative
missing-validation intervals.
There are $48\cdot4=192$ prespecified event bounds;
their union failure probability is at most
$192/7680=1/40$.

On the simultaneous coverage event, $p^-_{jk}\le p_{jk}$.
Because the thresholds are increasing and $\theta_0=0$,
the following pointwise inequality holds:
\[
 \sum_{k=1}^4(\theta_k-\theta_{k-1})
          \boldsymbol1_{\{G_j(x)\ge\theta_k\}}
 \le G_j(x).
\]
If the largest passed threshold is $\theta_s$, the
left side is exactly $\theta_s$; otherwise it is zero.
Taking expectations proves \eqref{eq:staircase120}.
Each individual lower bound $\theta_kp^-_{jk}$
is also valid.  The staircase statement uses all fixed
thresholds under the same simultaneous coverage event.

For the mean interval, a short proof of the needed
bounded-variable inequality is useful.  If $W\in[0,\tau]$,
the logarithmic moment generating function
$\psi(\lambda)=\log\mathbb E e^{\lambda W}$ has
$\psi''(\lambda)=\operatorname{Var}_\lambda(W)
\le\tau^2/4$ under its exponentially tilted law.
This variance bound follows from
$\mathbb E_\lambda(W-\tau/2)^2\le\tau^2/4$ and
the minimizing property of the mean.
Integrating twice from zero gives
\[
 \mathbb E e^{\lambda(W-\mathbb EW)}
 \le e^{\lambda^2\tau^2/8}.
\]
Independence, exponential Markov inequality, and minimization
over $\lambda$ yield for either tail
\[
 P\left(\left|\overline W-\mathbb EW\right|
       \text{ exceeds }r\text{ in that tail}\right)
 \le e^{-2nr^2/\tau^2}.
\]
Since $\overline\ell_j\le\overline G_j\le\overline u_j$,
the lower failure
$\mathbb EG_j<\overline\ell_j-r$ and upper failure
$\mathbb EG_j>\overline u_j+r$ are contained in the
corresponding two deviation events for the unobserved
true gain average.  Thus no numerical value error has
been ignored.

For the prespecified rational radius $r=33\tau/128$,
\[
 2nr^2/\tau^2=1089/128=:x.
\]
The strict inequality $e^x>3840$ can itself be checked
without a floating logarithm or exponential.
Write $S_m(z)=\sum_{k=0}^m z^k/k!$.
Since $x>17/2$, positivity of coefficients gives
\[
 S_{20}(x)\ge S_{12}(17/2)
       \ge S_{11}(8)(1+1/2)>2644(3/2)=3966>3840.
\]
The middle inequality follows by expanding
$(8+1/2)^k$: every term of
$S_{11}(8)S_1(1/2)$ occurs with nonnegative coefficient
in $S_{12}(8+1/2)$.
For the last lower bound, integer lower bounds for the
successive terms of $S_{11}(8)$ are
$1,8,32,85,170,273,364,416,416,369,295,215$;
their sum is $2644$, and several inequalities are strict.
Since $e^x\ge S_{20}(x)$, each one-sided mean failure
is strictly less than $1/3840$.
The $48$ two-sided mean intervals therefore fail with
total probability less than $96/3840=1/40$.

A final union bound over the event and mean families
gives failure probability at most $1/20$.
Intersecting valid bounds with the known population
range $[0,\tau]$ proves \eqref{eq:meangain120} and the
simultaneous $0.95$ assertion.
No independence between the two inference families is
used.  These are ordinary probability bounds
\citep{clopperpearson1934,hoeffding1963}, applied after
the original-objective safeguard makes their gain
range and events mathematically valid.
\end{proof}
